\chapter{Frontend}
\label{cap:frontend}

Next.js 16 con App Router, React 19 y TanStack Query. El cliente HTTP no se
escribe: se genera del OpenAPI de la API.

\section{Estructura de rutas}

Un solo grupo de rutas, \pkcode{(app)}, tres \emph{layouts} y una regla simple:
toda la superficie autenticada es cliente y consulta con TanStack Query; la
única página con obtención de datos en servidor es la vista pública compartida.

\begin{lstlisting}[language=bash]
app/
|- layout.tsx                    servidor  lang="es", fuentes, providers
|- page.tsx                      servidor  redirect("/collection")
|- sign-in/page.tsx              cliente
|- s/[token]/page.tsx            servidor  RSC asincrono, SIN autenticacion
|- api/auth/[...all]/route.ts    manejador de ruta -> Better Auth
`- (app)/                        la superficie autenticada
   |- layout.tsx                 cliente   el porton de sesion
   |- collection/[id|add]        cliente
   |- cards/[cardId]             cliente
   |- chat, scan, stats, compare cliente
   |- collectors/[id]            cliente
   |- trades/[new|publish]       cliente
   |- messages/[partnerId]       cliente   layout propio: la lista sobrevive
   `- notifications, profile     cliente
\end{lstlisting}

\begin{pkwarn}[No hay \pkcode{middleware.ts}]
La redirección a \pkcode{/sign-in} vive en un efecto de
\pkpath{app/(app)/layout.tsx}. Es decir: el HTML del grupo \pkcode{(app)} se
sirve a cualquiera, y lo que protege los \emph{datos} es que la API rechace toda
petición sin JWT válido, no el frontend. Es una decisión coherente con el modelo
—el frontend no es una frontera de confianza— pero conviene tenerla explícita.
\end{pkwarn}

La página pública \pkpath{app/s/[token]/page.tsx} es la excepción: hace
\pkcode{fetch} en servidor contra \pkcode{API\_INTERNAL\_URL} —la dirección del
contenedor, no la del navegador— con \pkcode{cache: "no-store"}, y llama a
\pkcode{notFound()} si la respuesta no es correcta.

\section{Autenticación en el cliente}

La cookie de sesión nunca sale del origen de Next. El recorrido completo:

\begin{enumerate}[leftmargin=1.4em,itemsep=2pt,topsep=3pt]
  \item El navegador guarda la cookie de Better Auth para \pkcode{localhost:3000}.
  \item \pkcode{getAccessToken()} llama a \pkcode{/api/auth/token} con
        \pkcode{credentials: "include"} —mismo origen, la cookie viaja— y recibe
        un JWT.
  \item Ese JWT se cachea en un módulo y se manda como
        \pkcode{Authorization: Bearer …} a FastAPI, en otro origen.
  \item FastAPI lo verifica contra el JWKS. Nunca ve la cookie.
\end{enumerate}

\begin{lstlisting}[language=Java]
let cachedToken: string | null = null;

export async function getAccessToken(): Promise<string | null> {
  if (cachedToken) return cachedToken;
  // Better Auth mints the JWT from the session cookie; the API never sees the cookie.
  const response = await fetch("/api/auth/token", { credentials: "include" });
  if (!response.ok) return null;
  const { token } = (await response.json()) as { token?: string };
  cachedToken = token ?? null;
  return cachedToken;
}
\end{lstlisting}

\pkcode{clearAccessToken()} se invoca al iniciar sesión, al cerrarla y al cambiar
de cuenta desde el perfil.

\section{El envoltorio del cliente generado}

Kubb genera funciones que aceptan un \pkcode{client}; la aplicación inyecta el
suyo, que hace tres cosas y ninguna más:

\begin{lstlisting}[language=Java]
const response = await client<TResponseData, TError, TRequestData>({
  ...config,
  headers: {
    ...(isJsonBody ? { "Content-Type": "application/json" } : {}),
    ...(config.headers as Record<string, string> | undefined),
    ...(token ? { Authorization: `Bearer ${token}` } : {}),
  },
});

if (response.status >= 400) {
  throw new ApiError(response.status, detailOf(response.data));
}
\end{lstlisting}

El orden de los tres \emph{spreads} es intencionado y está cubierto por pruebas:
\pkcode{Content-Type} va primero para que quien llama pueda sobrescribirlo,
\pkcode{Authorization} va último para que un encabezado obsoleto de quien llama
no pueda desplazar al token vivo, y un cuerpo \pkcode{FormData} se queda sin
\pkcode{Content-Type} a propósito, para que el navegador escriba el
\emph{boundary} del multipart.

\pkcode{ApiError} traduce las dos formas que FastAPI usa para \pkcode{detail}
—cadena, o lista cuyo primer elemento trae \pkcode{msg}— y cae a un mensaje en
español cuando el cuerpo no es ninguna de las dos.

\begin{pkwarn}[El envoltorio se inyecta por llamada]
\pkpath{kubb.config.ts} sólo elige \pkcode{fetch} sobre axios; el envoltorio se
pasa en cada sitio de llamada como \pkcode{\{ client: \{ client: apiClient \} \}}.
Aparece en 39 archivos. Olvidarlo no rompe la compilación: simplemente salta el
token \emph{y} el manejo de errores.
\end{pkwarn}

\section{Regeneración del cliente tipado}

\begin{lstlisting}[language=bash]
# 1. la API tiene que estar sirviendo, porque la entrada por omision es su URL
npm run dev            # o docker compose up -d api

# 2. desde apps/web
npx kubb generate      # o: OPENAPI_URL=./openapi.json npx kubb generate
\end{lstlisting}

\pkpath{kubb.config.ts} declara cinco plugins y cuatro destinos bajo
\pkpath{apps/web/lib/api/}: \pkpath{types/} (tipos TypeScript), \pkpath{zod/}
(esquemas de validación), \pkpath{clients/} (una función \pkcode{fetch} por
operación) y \pkpath{hooks/} (un \pkcode{useX} y un \pkcode{useXSuspense} por
operación). \pkcode{output.clean: true} borra el directorio antes de escribir,
por lo que cualquier retoque manual se pierde; cada archivo generado lleva
escrito \emph{«Do not edit manually»}.

\begin{pknote}[Detalles que ahorran una tarde]
No hay script de \pkcode{npm} para esto: la invocación es el binario de
\pkcode{@kubb/cli}, que sí está en \pkcode{devDependencies}. El árbol generado
\emph{se versiona} en git, y está excluido de las pruebas por
\pkpath{vitest.config.ts}. La opción \pkcode{extension: \{ ".ts": "" \}} existe
porque Kubb emite la extensión en las importaciones y TypeScript la rechaza sin
\pkcode{allowImportingTsExtensions}.
\end{pknote}

\section{TanStack Query}

\begin{lstlisting}[language=Java]
export function QueryProvider({ children }: { children: ReactNode }) {
  const [queryClient] = useState(() => new QueryClient());
  return <QueryClientProvider client={queryClient}>{children}</QueryClientProvider>;
}
\end{lstlisting}

Sin opciones. El comportamiento de producción es por tanto el de fábrica de
TanStack Query v5: \pkcode{staleTime: 0}, \pkcode{gcTime} de cinco minutos,
tres reintentos con retroceso exponencial y refetch al enfocar la ventana, al
montar y al reconectar. El inicializador perezoso del \pkcode{useState} es la
protección estándar del App Router contra compartir un cliente entre peticiones.

Que \pkcode{staleTime} sea cero es precisamente lo que hace importante a
\pkcode{CommittingRoute} (sección~\ref{cap:backend}): una mutación seguida de
invalidación provoca un refetch inmediato, y ese refetch no puede adelantarse al
commit.

\section{La URL como estado}

\pkpath{apps/web/lib/url-state.ts} evita deliberadamente \pkcode{useRouter} y
\pkcode{useSearchParams}:

\begin{quote}\small\itshape
\pkcode{router.replace} es una navegación aunque sólo se mueva un parámetro de
consulta: Next vuelve a pedir el payload de la ruta y el árbol se vuelve a
montar, así que cada clic en un filtro parpadeaba con un esqueleto y volvía a
pedir lo que ya estaba en pantalla. La History API escribe la misma URL sin ida
y vuelta\dots{} El precio es que el estado ahora es de React, así que todo lo que
lea la URL tiene que pasar por este hook.
\end{quote}

\par\addvspace{8pt}\noindent\begin{pkbox}{colspec={X[1.2,l]X[3.4,l]}}
  Parte de la API & Comportamiento \\
  \pkcode{useUrlState(initial)} & Devuelve \pkcode{[params, set]}. \pkcode{params} es un \pkcode{URLSearchParams} sembrado desde \pkcode{window.location.search} (a prueba de SSR) y completado con \pkcode{initial} \emph{sólo} en las claves que la URL no traiga. \\
  \pkcode{set(next)} & Fusiona; \pkcode{undefined} o \pkcode{""} borran la clave; escribe con \pkcode{history.replaceState}, sin entrada de historial y sin navegación. \\
  \pkcode{popstate} & Un escucha relee la URL, así que el botón «atrás» funciona; se retira al desmontar. \\
\end{pkbox}\par\addvspace{8pt}

Lo usan diez pantallas. El patrón habitual es fijar un filtro y limpiar la
página en la misma escritura: \pkcode{setParam(\{ q: next, p: undefined \})}.

\section{El sistema de diseño}

\pkpath{packages/ui} son 25 componentes shadcn construidos sobre \textbf{Base
UI}, no sobre Radix, más un archivo de \emph{tokens}. Tailwind~4 sin archivo de
configuración: todo vive en CSS.

\par\addvspace{8pt}\noindent\begin{pkbox}{colspec={X[1.3,l]X[3.4,l]}}
  Decisión & Razón, tal como está comentada en el CSS \\
  \pkcode{--primary: oklch(0.56 0.23 27)} & El rojo de la pokédex. Es el único color saturado del sistema fuera del arte de las cartas. \\
  \pkcode{--ring: oklch(0.58 0.15 250)} & Azul \emph{a propósito}: «primario y destructivo son ambos rojos, así que un anillo de foco rojo sobre un campo es indistinguible de que ese campo sea inválido». \\
  18 tokens \pkcode{--type-*} & Un color por tipo de Pokémon, re-exportados al tema de Tailwind. \\
  \pkcode{TYPE\_VAR} en TS & Mapea tipo → \emph{nombre de variable} y no a una clase: «un color de tipo se usa como relleno, resplandor, anillo y parada de degradado, y sólo una variable cruda sirve para las cuatro». \\
  \pkcode{@utility slab / glass / aura / foil} & Utilidades propias de Tailwind~4 para la superficie, la cromática flotante, el resplandor teñido por tipo y el holográfico. \\
  \pkcode{forcedTheme="light"} & La aplicación es clara y sólo clara: un tema por omisión seguiría al sistema operativo y dejaría a quien lo tenga oscuro en un tema que ya no se mantiene. \\
  \pkcode{prefers-reduced-motion} & Todas las animaciones y transiciones se recortan a 0.01~ms. \\
\end{pkbox}\par\addvspace{8pt}

El paquete no se compila: exporta un archivo por subruta
(\pkcode{@workspace/ui/components/button}) y \pkpath{apps/web/tsconfig.json}
apunta el alias directamente a \pkpath{packages/ui/src}, así que no existe
artefacto intermedio que pueda quedar viejo.

\section{Empaquetado}

\pkcode{output: "standalone"} traza el grafo de dependencias en un paquete
autocontenido, de modo que la imagen de ejecución no lleva \pkpath{node\_modules}.
Esa elección obliga a copiar a mano \pkpath{.next/static} y \pkpath{public} en la
etapa final, y a instalar aparte las dependencias de \pkpath{apps/web/scripts},
que es un paquete npm propio: Next empaqueta \pkcode{better-auth} dentro de sus
\emph{chunks} de servidor en vez de copiar el paquete, así que la migración no
puede tomarlo prestado del árbol trazado.

\begin{pkwarn}[\pkcode{NEXT\_PUBLIC\_API\_URL} es argumento de construcción, no variable de ejecución]
Se hornea en el \emph{bundle} del cliente al compilar, así que tiene que ser la
URL que alcanza un navegador del anfitrión, no la que usa el contenedor.
Cambiarla exige reconstruir la imagen, no reiniciarla. Su contraparte
\pkcode{API\_INTERNAL\_URL} sí es una variable de ejecución normal, y sólo la lee
el servidor en la página pública compartida.
\end{pkwarn}
